\section*{Disponibilidad del código y los artefactos}
\phantomsection
\label{sec:disponibilidad-codigo}

El código fuente, las pruebas automatizadas, los scripts de ejecución,
el notebook y los artefactos experimentales asociados con este trabajo
se encuentran disponibles públicamente en el repositorio:

\begin{center}
\url{https://github.com/hsancheza90-art/keccak-milp-criptografia}
\end{center}

Para garantizar la trazabilidad de los resultados, la versión utilizada
en este artículo se identifica mediante la etiqueta:

\begin{center}
\path{keccak-milp-experiments-v1}
\end{center}

La etiqueta corresponde al commit:

\begin{center}
\path{0b089aa}
\end{center}

La rama \path{feature/chi-milp} contiene el desarrollo del modelo y
puede recibir modificaciones posteriores. Por esta razón, la etiqueta
estable y el identificador del commit constituyen las referencias
recomendadas para reproducir los resultados presentados.

\subsection*{Estructura de los recursos publicados}

La Tabla~\ref{tab:published-resources} resume los recursos principales
del repositorio.

\begin{table}[H]
\centering
\caption{Recursos disponibles para reproducción y auditoría.}
\label{tab:published-resources}
\begin{tabularx}{\textwidth}{
    >{\raggedright\arraybackslash}p{0.40\textwidth}
    X
}
\toprule
\textbf{Recurso} & \textbf{Contenido} \\
\midrule

\path{src/keccak_milp/}
&
Implementación de referencia de Keccak reducido, formulación MILP,
modelo diferencial pareado y funciones de búsqueda.
\\

\path{scripts/run_active_sbox_experiments.py}
&
Punto de entrada para ejecutar la matriz de experimentos en los modos
\path{quick} y \path{full}.
\\

\path{tests/}
&
Pruebas unitarias, de integración, reconstrucción de testigos y
validación de artefactos.
\\

\path{notebooks/keccak_milp_active_sboxes.ipynb}
&
Notebook académico ejecutado, con tablas, figura, resultados,
interpretación y pruebas dirigidas.
\\

\path{outputs/results/active_sbox_results.csv}
&
Resultados consolidados de los seis experimentos.
\\

\path{outputs/results/active_sbox_witnesses.json}
&
Estados, soportes y datos necesarios para reconstruir los testigos.
\\

\path{outputs/results/experiment_environment.json}
&
Información del entorno y de la ejecución experimental.
\\

\path{outputs/figures/active_sbox_bounds.png}
&
Figura comparativa de las cotas inferiores y superiores.
\\

\path{report/}
&
Fuentes en LaTeX y versión compilable del artículo académico.
\\

\bottomrule
\end{tabularx}
\end{table}

\subsection*{Reproducción mínima}

Una reproducción funcional de los resultados permanentes puede
realizarse desde la raíz del repositorio mediante:

\begin{lstlisting}[language=bash]
python -m pytest -q --disable-warnings
python scripts/run_active_sbox_experiments.py --mode quick
\end{lstlisting}

Para volver a ejecutar las búsquedas exhaustivas restringidas de tres
rondas debe utilizarse:

\begin{lstlisting}[language=bash]
python scripts/run_active_sbox_experiments.py --mode full
\end{lstlisting}

El informe se compila desde la carpeta \path{report/}:

\begin{lstlisting}[language=bash]
latexmk -pdf -interaction=nonstopmode -halt-on-error main.tex
\end{lstlisting}

\subsection*{Referencia de la versión experimental}

La combinación de los siguientes elementos identifica de forma
unívoca la versión del trabajo utilizada en el artículo:

\begin{itemize}
    \item repositorio:
    \url{https://github.com/hsancheza90-art/keccak-milp-criptografia};

    \item etiqueta:
    \path{keccak-milp-experiments-v1};

    \item commit:
    \path{0b089aa};

    \item ejecutor:
    \path{scripts/run_active_sbox_experiments.py};

    \item notebook:
    \path{notebooks/keccak_milp_active_sboxes.ipynb}.
\end{itemize}

Esta información permite recuperar el código, verificar los testigos y
reconstruir las tablas y figuras sin depender de la evolución futura de
la rama de desarrollo.
